% !TeX root = ../thesis.tex
% ===================== ВВЕДЕНИЕ =====================
\chapter*{Введение}
\addcontentsline{toc}{chapter}{Введение}

\textbf{Актуальность.}
Описание течений разреженного газа, в которых длина свободного пробега молекул сравнима с характерным размером задачи
(число Кнудсена $Kn\sim 1$), выходит за рамки уравнений сплошной среды и требует кинетического подхода~--- решения
уравнения Больцмана.
Такие режимы возникают в гиперзвуковой аэродинамике (обтекание спускаемых аппаратов в верхних слоях атмосферы), в
вакуумной и микроканальной технике, при разделении изотопов, в физике плазмы и при моделировании ударных волн.
Главная трудность~--- нелинейный интеграл столкновений высокой кратности в правой части, поэтому разработка эффективных
и \emph{консервативных} солверов кинетического уравнения, пригодных для счёта на суперкомпьютерных системах, остаётся
востребованной задачей вычислительной физики.

\textbf{Обзор подходов.}
Наиболее распространён метод прямого статистического моделирования (ПСМ, DSMC) Г.~Бёрда~\cite{Bird}, однако вблизи
равновесия он страдает от статистического шума, требующего большого числа частиц.
Детерминированные методы дискретных ординат (Бродвелл, Кабаннес) разрешают лишь те столкновения, которые точно попадают
в узлы скоростной сетки, что сужает класс учитываемых взаимодействий.
Интерполяционные методы (В.\,В.~Аристов и др.)~\cite{Aristov} разрешают все столкновения, но при наивной интерполяции
нарушают сохранение импульса и энергии, порождая «паразитные» источники.
\emph{Консервативный проекционный метод} Ф.\,Г.~Черемисина~\cite{Tcher,TcherJCP} устраняет эти недостатки: отклик
интеграла столкновений проецируется на узлы скоростной сетки так, что масса, импульс и энергия сохраняются на уровне
разностной схемы, а логарифмическая (геометрическая) интерполяция гарантирует положительность функции распределения и
выполнение $H$-теоремы.
Для вычисления многомерного интеграла столкновений вместо метода Монте-Карло применяются теоретико-числовые сетки
Коробова~\cite{Korobov} с погрешностью $O\!\left((\ln p)^{8}/p\right)$ (для 8-мерной сетки) против
$O\!\left(1/\sqrt{N}\right)$ у метода Монте-Карло (здесь $p$~--- число узлов сетки Коробова, $N$~--- число случайных
выборок).

\textbf{Предмет, цель и задачи.}

\emph{Объект исследования}~--- уравнение Больцмана для функции распределения $f(t,x,\xv)$.

\emph{Предмет}~--- численная схема его решения консервативным проекционным методом и её применение к задачам релаксации
и переноса.

\emph{Цель работы}~--- разработать и верифицировать программный солвер, строго сохраняющий макропараметры, и применить
его к задаче переноса (ударная волна перед поршнем) с проведением расчётов на суперкомпьютерном кластере.

\emph{Задачи}
\begin{enumerate}
  \item Реализовать общую численную схему проекционного метода на сетках Коробова;
  \item верифицировать солвер на пространственно-однородных задачах релаксации (квази-шок и бинарная смесь);
  \item решить неоднородную задачу о поршне со схемой расщепления по физическим процессам;
  \item проконтролировать законы сохранения и $H$-теорему в дискретном пространстве скоростей;
  \item выполнить расчёты на суперкомпьютерном кластере и оценить путь к промышленному солверу.
\end{enumerate}

\textbf{Инструменты.}

Уравнение Больцмана в модели твёрдых сфер;

консервативный проекционный метод Черемисина;

8-мерные теоретико-числовые сетки Коробова;

схема расщепления Странга и TVD-схема переноса;

кластер ЦКП НИЦ «Курчатовский институт».

\textbf{Личный вклад автора.}
На основе известного проекционного метода Черемисина и базового ядра вычисления интеграла столкновений автором
самостоятельно реализованы:

неоднородный солвер задачи о поршне (TVD-схема переноса в потоковой форме, граничные условия через слои фиктивных ячеек,
расщепление Странга);

адаптация оператора столкновений к одномерной по пространству задаче (4-кратная поперечная симметрия вместо 8-кратной,
контроль положительности по реальному приращению);

векторизация по пространственным ячейкам с предвычислением геометрии столкновений (ускорение примерно на два порядка).

Автором выполнены также постановка всех расчётов на кластере, диагностика кинетического предвестника и подбор размера
области, количественная проверка соотношений Ренкина--Гюгонио.
Готовыми взяты сам проекционный метод и теория Черемисина, оптимальные коэффициенты сеток Коробова и базовая реализация
ядра интеграла столкновений из задачи релаксации.

\textbf{Научная и практическая значимость.}
С научной точки зрения работа подтверждает, что консервативный проекционный метод с теоретико-числовыми сетками Коробова
обеспечивает строгое выполнение законов сохранения и $H$-теоремы как в однородных задачах релаксации, так и в
неоднородной задаче переноса.
Практическая значимость состоит в том, что разработанный солвер пригоден в качестве референсного инструмента для
верификации более быстрых (в том числе статистических) методов, а его алгоритмическая структура естественно переносится
на суперкомпьютерные системы.
Полученные результаты могут быть использованы при моделировании ударно-волновых течений разреженного газа и процессов
установления равновесия в газовых смесях.

\textbf{Апробация.}
Результаты доложены на 68-й Всероссийской научной конференции МФТИ в честь 130-летия со дня рождения Н.\,Н.~Семёнова
(апрель 2026~года), где работа отмечена призовым местом; тезисы приняты к публикации.

\textbf{Структура работы.}

В главе~1 изложены математическая модель и численный метод.

В главе~2 солвер верифицирован на двух задачах релаксации.

В главе~3 решена задача переноса о поршне.

В главе~4 описаны программная реализация и расчёты.

В заключении сформулированы выводы и план будущих исследований.

